problem:  
Let $\mathcal{M}$ denote the space of smooth, simple, closed embeddings of $S^{1}$ into $\mathbb{R}^{2}$, modulo $\mathrm{Diff}(S^{1})$. Endow $\mathcal{M}$ with the first‑order Sobolev (elastic) Riemannian metric
\[
G_c(h,k) = \int_{S^{1}} \langle (1 - A\partial_{\theta}^{2})\,h , k \rangle\,|c'(\theta)|\,d\theta,
\]
where $h,k$ are normal vector fields along $c$ and $A>0$ is fixed.  
This induces a non‑degenerate geodesic distance $d_G$ on $\mathcal{M}$.  

Define $\overline{\mathcal{M}}$ to be the metric completion of $(\mathcal{M},d_G)$, and let $\mathrm{Imm}^{1,2}_{\mathrm{fin}}(S^{1},\mathbb{R}^{2})$ denote the space of $H^{1}$‑immersions $c$ with $|c'|>0$ a.e., finite total curvature in the sense that the curvature vector $\kappa_c n_c$ defines a finite Radon measure, and at most finitely many transversal self‑intersections.

**Question (in two parts):**  
(a) **Characterization:** Is every $d_G$‑Cauchy sequence of smooth embeddings (modulo reparametrization) represented by an element of $\mathrm{Imm}^{1,2}_{\mathrm{fin}}$?  
(b) **Realization:** Conversely, does every $c \in \mathrm{Imm}^{1,2}_{\mathrm{fin}}$ arise as the $d_G$‑limit of some sequence of smooth embeddings (modulo $\mathrm{Diff}(S^1)$)?  

A positive pair of answers would identify $\overline{\mathcal{M}}$ as being isometric to the quotient of $\mathrm{Imm}^{1,2}_{\mathrm{fin}}(S^{1},\mathbb{R}^{2})$ by $\mathrm{Diff}(S^1)$ and would yield a concrete geometric stratification of the completion according to the number and configuration of self‑intersections.

Why is it a "good" problem:  
This version distills a major open question in infinite‑dimensional Riemannian geometry into two precise and logically independent subquestions about metric completion and weak limits under Sobolev metrics.  

1. **Profound insight:** The problem addresses the boundary between smooth and weak geometric objects—what shapes emerge as metric limits when smooth embeddings degenerate under a Sobolev energy that controls both position and curvature. Understanding this passage probes the analytic structure underlying shape geometry and extends classical regularity theory into infinite-dimensional settings.  

2. **Bridge between fields:** The question blends infinite‑dimensional Riemannian geometry, geometric measure theory, and analysis of Sobolev immersions. Establishing such a correspondence would conceptually link the metric completion problem in geometric analysis to the regularity and singularity questions studied in curvature‑bounded geometric measure theory.  

3. **Extreme simplicity:** The problem can be summarized in an intuitive sentence: *“What are the limit shapes when smooth plane curves degenerate under the Sobolev elastic metric?”* Despite its brevity and conceptual accessibility, it demands major advances in nonlinear analysis, reparametrization‑invariant metric geometry, and measure‑theoretic curvature theory—representing a quintessential “narrow entrance, profound depth” research problem.